\documentclass[a4paper, 12pt]{memoir} %report
\usepackage[utf8]{inputenc}
\usepackage[spanish]{babel} % dictionary
\usepackage{subfiles} % organize code in different files
\usepackage[a4paper, margin=2.5cm]{geometry} % page size and margins
\usepackage{float} % positioning fibures and labels
\usepackage{threeparttablex} % fancy tables with three parts, top, mid and bottom
\usepackage{booktabs, caption, makecell}
\usepackage{graphicx} % add figures
\usepackage{subcaption} % add two subfigures at the same figure
\graphicspath{ {figures/} } % path to figures
\usepackage{amsmath} % math in latex
\usepackage{marvosym} % euros symbol
\usepackage{multirow} % multirow and multicolumn cells in tables
\usepackage{longtable} % for 'longtable' environment
\usepackage{pdflscape} % for 'landscape' environment
\usepackage{hyperref} % add hyperlinks in the index
\usepackage[acronym, toc, nonumberlist]{glossaries} % glossary and acronyms (toc to show in contents)
\usepackage{libertine} % fancy title
\usepackage{afterpage} % set big tables position

\usepackage{fancyhdr} %for headers and footers
\pagestyle{fancy}
\fancyhf{}

\fancyhead[LE,RO]{\thepage} %for having the page number left in even and right in odd pages
\fancyhead[RE]{\leftmark}
\fancyhead[LO]{ \nouppercase{\rightmark}}

\providecommand{\keywords}[1]{\textbf{\textit{Palabras clave---}} #1} % Defining keywords

%\renewcommand\thetable{\Alph{Tabla}}
%\renewcommand{\thetable}{\Roman{table}}
%\renewcommand\thefigure{S\roman{figure}}

\renewcommand*\entryname{Name}
\renewcommand*\descriptionname{Prueba}
\setglossarystyle{superheader}

\renewcommand*{\glsnamefont}[1]{\bfseries #1}

\title{TFG}
\author{pedro.arias }
\date{September 2019}

\makeglossaries

\newglossaryentry{Pi}
{
  name={\ensuremath{P_i}},
  description={Punto de referencia},
  sort=1
}

\newglossaryentry{coord}
{
  name={\ensuremath{x_i, y_i}},
  description={Coordenadas de un punto de referencia},
  sort=2
}

\newglossaryentry{S}
{
  name={\ensuremath{S}},
  description={Punto solución},
  sort=3
}

\newglossaryentry{coord_sol}
{
  name={\ensuremath{\hat{x}, \hat{y}}},
  description={Coordenadas de punto solución},
  sort=4
}

\newglossaryentry{theta}
{
  name={\ensuremath{\theta_i}},
  description={Ángulo de llegada},
  sort=5
}

\newglossaryentry{r}
{
  name={\ensuremath{r_i}},
  description={Rango, distancia entre dos puntos},
  sort=6
}

\newglossaryentry{c}
{
  name={\ensuremath{c}},
  description={Velocidad de la luz en el vacío},
  sort=7
}

\newglossaryentry{t_i}
{
  name={\ensuremath{t_i}},
  description={Tiempo de llegada, ToA},
  sort=8
}

\newglossaryentry{t}
{
  name={\ensuremath{t}},
  description={Tiempo},
  sort=9
}

\newglossaryentry{N}
{
  name={\ensuremath{N}},
  description={Número de referencias conocidas},
  sort=10
}

\newglossaryentry{p_r}
{
  name={\ensuremath{P_R}},
  description={Potencia o intensidad de señal recibida, RSSI},
  sort=11
}

\newglossaryentry{p_t}
{
  name={\ensuremath{P_T}},
  description={Potencia de señal transmitida},
  sort=12
}

\newglossaryentry{g_t}
{
  name={\ensuremath{G_T}},
  description={Ganancia de señal transmitida},
  sort=13
}

\newglossaryentry{g_r}
{
  name={\ensuremath{G_R}},
  description={Ganancia de señal recibida},
  sort=14
}

\newglossaryentry{lambda}
{
  name={\ensuremath{\lambda}},
  description={Longitud de onda de la señal},
  sort=15
}

\newglossaryentry{d}
{
  name={\ensuremath{d}},
  description={Distancia},
  sort=16
}

\newglossaryentry{n}
{
  name={\ensuremath{n}},
  description={Constante de propagación de Friis},
  sort=17
}

\newglossaryentry{coef}
{
  name={\ensuremath{\alpha, \beta, \gamma, \delta, \mu, \nu, \omega}},
  description={Coeficientes del modelo propuesto},
  sort=18
}

\newglossaryentry{carga de pago}
{
  name={\ensuremath{CP}},
  description={Carga de pago},
  sort=19
}

\newglossaryentry{infra_mant}
{
  name={\ensuremath{IM}},
  description={Costes de infraestructura y mantenimiento},
  sort=20
}

\newglossaryentry{usuario_final}
{
  name={\ensuremath{UF}},
  description={Costes de usuario final},
  sort=21
}

\newglossaryentry{rango}
{
  name={\ensuremath{RAN}},
  description={Rango, cobertura},
  sort=22
}

\newglossaryentry{pos}
{
  name={\ensuremath{POS}},
  description={Posicionamiento},
  sort=23
}

\newglossaryentry{act}
{
  name={\ensuremath{ACT}},
  description={Actitud, orientación},
  sort=24
}

\newglossaryentry{obs}
{
  name={\ensuremath{OBS}},
  description={Obstáculos},
  sort=25
}


% --------------------------------------------------------------%

\newacronym{ips}{IPS}{Indoor Positioning System}

\newacronym{gnss}{GNSS}{Global Navigation Satellite System}

\newacronym{gps}{NAVSTAR-GPS}{Navigation System Time And Ranging-Global Positioning System}

\newacronym{glonass}{GLONASS}{Globalnaja Nawigazionnaja Sputnikowaja Sistema}

\newacronym{vcl}{VCL}{Visible Ligth Communication}

\newacronym{ir}{IR}{Infrared}

\newacronym{rf}{RF}{Radio Frecuency}

\newacronym{uwb}{UWB}{Ultra-wideband}

\newacronym{wsn}{WSN}{Wireless Sensor Network}

\newacronym{wlan}{WLAN}{Wireless Local Area Network}

\newacronym{wifi}{Wi-Fi}{Wireless Fidelity}

\newacronym{rfid}{RFID}{Radio Frecuency Identification}

\newacronym{nfc}{NFC}{Near Field Communication}

\newacronym{toa}{ToA}{Time of Arrival}

\newacronym{tof}{ToF}{Time of Flight}

\newacronym{aoa}{AoA}{Angle of Arrival}

\newacronym{tdoa}{TDoA}{Time Difference of Arrival}

\newacronym{rssi}{RSSI}{Received Signal Strength Indicator}

\newacronym{ins}{INS}{Inertial Navigation System}

\newacronym{imu}{IMU}{Inertial Measurement Unit}

\newacronym{2d}{2D}{Two Dimensions}

\newacronym{3d}{3D}{Three Dimensions}

\newacronym{pl}{PL}{Path Loss}

\newacronym{k-nn}{k-NN}{k-Nearest Neighbor}

\newacronym{svm}{SVM}{Support Vector Machine}

\newacronym{smp}{SMP}{Smallest M-vertex Polygon}

\newacronym{slam}{SLAM}{Simultaneous Location and Mapping}

\newacronym{los}{LoS}{Line of Sight}

\newacronym{led}{LED}{Light Emitting Diode}

\newacronym{vo}{VO}{Visual Odometry}

\newacronym{lidar}{LiDAR}{Light Detection and Ranging / Laser Imaging Detection and Ranging}

\newacronym{uav}{UAV}{Unmanned Aerial Vehicle}

\newacronym{lan}{LAN}{Local Area Network}

\newacronym{ssid}{SSID}{Service Set Identifier}

\newacronym{ble}{BLE}{Bluetooth Low Energy}

\newacronym{gnc}{GNC}{Guidance, Navigation and Control}

\newacronym{mav}{MAV}{Micro Aerial Vehicle}

\newacronym{urjc}{URJC}{Universidad Rey Juan Carlos}

\newacronym{ieee}{IEEE}{Institute of Electrical and Electronics Engineers}

\newacronym{ekf}{EKF}{Extended Kalman Filter}

\newacronym{epm}{EPM}{Ellipse Propagation Model}

\newacronym{fis}{FIS}{Fuzzy Inference System}

\newacronym{rtof}{RToF}{Roundtrip Time of Flight}

\newacronym{poa}{PoA}{Phase of Arrival}

\newacronym{pdoa}{PDoA}{Phase Difference of Arrival}

\newacronym{dolphin}{DOLPHIN}{Distributed Object Locating System for Physical-space Internetworking}

\newacronym{blpa}{BLPA}{Bluetooh Local Positioning Application}

\newacronym{landmarc}{LANDMARC}{Location Identification based on Dynamic Active RFID Calibration}

\newacronym{rtls}{RTLS}{Real Time Location System}

\newacronym{bds}{BDS}{BeiDou System}

\newacronym{loran}{LORAN}{LOng RAnge Navigation}

\newacronym{irnss}{IRNSS}{Indian Regional Navigation Satellite System}

\newacronym{qzss}{QZSS}{Quasi-Zenith Satellite System}

\newacronym{kf}{KF}{Kalman Filter}

\newacronym{ukf}{UKF}{Unscented Kalman Filter}

\newacronym{pid}{PID}{Proportional–Integral–Derivative}

\newacronym{mpc}{MPC}{Model Predictive Control} % Acronyms and symbols

\begin{document}

\renewcommand*\listtablename{Índice de tablas}
\renewcommand\spanishtablename{Tabla}
\renewcommand\theadfont{\bfseries}


\frontmatter

% \maketitle
\begin{figure}[t]
    \centering\includegraphics[width=0.5\textwidth]{photos/urjc_logo_letters.jpg}
\end{figure}

\begin{center}
    \textsc{ \Large{ESCUELA TÉCNICA SUPERIOR DE INGENIERÍA DE TELECOMUNICACIÓN \\}}
    \vspace{10mm}
	\textnormal{ \Large{GRADO EN INGENIERÍA AEROESPACIAL EN AERONAVEGACIÓN\\}}
	\vspace{25mm}
    \textup{\LARGE{\textbf{TRABAJO FIN DE GRADO}}} \\
    
    \vspace{25mm}
    \textnormal{ \Large{SISTEMAS DE NAVEGACIÓN EN INTERIORES APLICADO A MICRO-VEHÍCULOS AÉREOS (MAVs): APROXIMACIÓN A UN DISEÑO CONCEPTUAL}}
\end{center}

\vspace{30mm}

\begin{center}
    \textnormal{\normalsize{\textbf{Autor:} Pedro Arias Pérez \\}}
    \textnormal{\normalsize{\textbf{Co-tutor:} Raúl de Celis Fernández \\}}
    \textnormal{\normalsize{\textbf{Co-tutor:} Hodei Urrutxua Cereijo \\}}

\end{center}

\vspace{20mm}

\begin{center}
   \normalsize{Curso académico 2019/2020} 
\end{center}

\cleardoublepage

\begin{vplace}[0.7]
    \begin{flushright}
	    \textit{Aos meus pais,\\ que dedicaron a súa vida á miña educación.}
    \end{flushright}
\end{vplace}
\cleardoublepage

\addcontentsline{toc}{chapter}{Resumen}
\subfile{chapters/abstrac}

\chapter*{Agradecimientos}
\addcontentsline{toc}{chapter}{Agradecimientos}
\subfile{chapters/acknowledgments}
\clearpage

% Indice
\tableofcontents
\clearpage

% Indices de figuras y tablas
\listoffigures
\clearpage
\listoftables

% Acronimos y simbolos
\printglossary[type=\acronymtype, title=Acrónimos, toctitle=Acrónimos, style=superheader] %alttreehypergroup
\cleardoublepage
\printglossary[title=Símbolos, toctitle=Símbolos, style=superheader]
\cleardoublepage

\mainmatter

\chapter{Introducción} \label{chapter:introduccion}
\subfile{chapters/01-introduccion}
\clearpage

\chapter{Técnicas de posicionamiento en interiores} \label{chapter:tecnicas}
\subfile{chapters/02-tecnicas}

\chapter{Tecnologías para posicionamiento en interiores} \label{chapter:tecnologias}
\subfile{chapters/03-tecnologias}
\clearpage

\chapter{Comparativa de sistemas} \label{chapter:comparativa}
\subfile{chapters/04-comparacion}
\clearpage

\chapter{Caso de estudio} \label{chapter:estudio}
\subfile{chapters/05-casoestudio}
\clearpage

\chapter{Conclusiones y líneas futuras} \label{chapter:conclusiones}
\subfile{chapters/06-conclusiones}

\glsaddall % show all acronyms and symbols

% Referencias
\bibliographystyle{plain} % named or plain
\bibliography{refs} % "refs.bib" file

\end{document}


%%%% TO DO %%%%
% 2. Abstrac en inglés -> NO
% 7. Cambiar intro ->  remontarme a la necesidad de posicionar -> NO

% 6. Unificar estilo referencias. REF cambiar por orden de aparicion

% 1. Erratas menores
% 2. Objetivos secundarios
% 3. Mencionar que no esta claro que la hibridacion sea la mejor solucion
% 5. Revisar tablas 4.1 -> Magnetica
% 6. extender 1.1 con citas

%%%%%%%%%%%%%%%%%%%%%%%%%%%%%%%%%%%%%%%%%%%%%%%%%%%%%%%%%%

%%%% ESTRUCTURA %%%%
% 1. Intro
% 2. Estudio
% 3. Analisis solucion
% 4. Modelado
% 5. Conclusión